\section{The estimated gain and the remaining upper-source measures}
\label{sec:v93-upper}
The positive-growth budget can be estimated on its entire noncritical
part, uniformly in every previously fixed positive cutoff. This leaves
actual ancestral entrances and critical positive growth, rather than
an integrated oscillation norm, in the pointwise inequality.

For a fixed gradient threshold $\rho>0$, let
$G_{c,r}^{\varepsilon,<\rho}$ and
$G_{c,r}^{\varepsilon,\ge\rho}$ be the terms in
\eqref{eq:v93-ancestral-budget} restricted at the intermediate roof
to $z<\rho$ and $z\ge\rho$, respectively. Define their sharp versions
by removing the ramp. The cutoff is on the physical intermediate
state, not on an initially transported surrogate source.

\begin{theorem}[Complete noncritical positive gain]
\label{thm:v93-gain}
At every finite count and almost every roof,
\begin{equation}\label{eq:v93-gain-finite}
 G_{c,r}^{\varepsilon,\ge\rho,\sharp}(t)
 \le C(1+\varepsilon^{-1})\int_{t-r}^t v^\rho(u)\,du,
\end{equation}
where $v^\rho$ is the actual truncated first-inverse-gradient source
in Input~\ref{input:v93-analytic}. Consequently
\begin{equation}\label{eq:v93-gain-ordered}
 \limsup_m\sup_{R,n,k,c}m^2
       \|G_{c,r}^{\varepsilon,\ge\rho,\sharp}\|_\infty
        \le C(1+\varepsilon^{-1})r.
\end{equation}
The supremum over $c$ ranges over all Borel prefixes at that count.
The leading constant is independent of $\rho$, $r$ and the two
original exceptional-source cutoffs. Each $\rho,r>0$ is fixed before
the collision limit. The theorem does not assert its conclusion at
$\rho=0$.
\end{theorem}
\begin{proof}
Before applying the absolute bound of Theorem~\ref{thm:v93-traces},
keep the physical density in the product rule:
\[
 [(f_\sigma^\varepsilon)'(u)]_+
 \le C(1+\varepsilon^{-1})
                  A_\sigma(u)I_m(\Psi_\sigma(u))/z(\Psi_\sigma(u)).
\]
For each intermediate roof $u$, ancestors of a fixed selected target
form a subfamily of the original exact fiber, counted once by the
Borel segment map. Integrating this density over that subfamily and
then over $u\in(t-r,t)$ is bounded by the right side of
\eqref{eq:v93-gain-finite}. No trace is equated to unweighted flux,
and no factor counting physical words is inserted. Apply the
positive-cutoff estimate \eqref{eq:v93-cutoff-input} with interval
length $r$ to prove \eqref{eq:v93-gain-ordered}. Its finite-count
remainder depends on $\rho$; the legal ordered estimate has one
leading constant for all previously fixed positive $\rho$.
\end{proof}

Combining Theorems~\ref{thm:v93-pointwise} and \ref{thm:v93-gain}
gives the following estimate on the original, unrenormalized source:
\begin{equation}\label{eq:v93-upper-finite}
 R_c^\varepsilon(t)\le
 \frac1r\int_{t-r}^tR_c^\varepsilon(s)\,ds
 +B_{c,r}^\varepsilon(t)+G_{c,r}^{\varepsilon,<\rho}(t)
 +C(1+\varepsilon^{-1})\int_{t-r}^t v^\rho(u)\,du.
\end{equation}
Every term is formed at finite count. In particular the left side
includes selected points with arbitrarily small initial gradient and
arbitrarily short total residence.

\subsection{A source-specific upper-spike certificate}
For the sharp budgets put
\[
 H_{c,r,\rho}^\varepsilon(t)
 =B_{c,r}^{\varepsilon,\sharp}(t)
       +G_{c,r}^{\varepsilon,<\rho,\sharp}(t).
\]
Fix a roof $t$ and a level $\lambda>0$. If
$m^2R_c^\varepsilon(t)\ge\lambda$ and
\begin{equation}\label{eq:v93-spike-test}
 m^2H_{c,r,\rho}^\varepsilon(t)
 +Cm^2(1+\varepsilon^{-1})\int_{t-r}^t v^\rho\le\lambda/2,
\end{equation}
then \eqref{eq:v93-persistence} proves
\begin{equation}\label{eq:v93-spike-width}
 m^2R_c^\varepsilon(s)\ge\lambda/2
            \quad\text{for almost every }s\in(t-r,t).
\end{equation}
This is persistence on an entire prescribed interval, not only a
weak bound for the measure of some superlevel set. It is quantitative
and holds for the original selector, but condition
\eqref{eq:v93-spike-test} is not asserted uniformly at every roof.

For example, if a fixed interval containing $(t-r,t)$ has
$\int m^2R_c^\varepsilon<\lambda r/2$, then an upper spike forces
the left side of \eqref{eq:v93-spike-test} to exceed $\lambda/2$.
Alternatively its critical measurable-set norm must be at least
$(\lambda/2)r^{31/32}$ unless the same physical budget is large.
The already proved source-mass and critical-norm decay therefore
force any surviving high spike into a concentration of actual
ancestral entrance traces or critical positive optical growth.
The noncritical gain has already been paid by
\eqref{eq:v93-gain-ordered}.

\subsection{An upper error estimate in the original topology}
For the full bad source take $c_\sigma=b_\sigma$ and omit $c$.
On the original central target define
\begin{align}
 \mathfrak I_{\varepsilon,M}(r)
   &=\limsup_m\sup_{R,n,k}m^2\esssup_{t\in\mathcal C(M)}
                                        B_r^\varepsilon(t),\notag\\
 \mathfrak C_{\varepsilon,M}(r)
   &=\lim_{\rho\downarrow0}\limsup_m\sup_{R,n,k}m^2
        \esssup_{t\in\mathcal C(M)}G_r^{\varepsilon,<\rho}(t).
                                             \label{eq:v93-forcing-charges}
\end{align}
The latter limit is monotone and extended nonnegative. It is taken
\emph{after} the normalized essential supremum and collision limit.
These positive quantities are not signed cancellations of exits.

\begin{theorem}[The scalar upper half with its positive mechanisms separated]
\label{thm:v93-upper-endpoint}
Under the retained analytic inputs, for all large fixed $B$ and fixed
$r>0$, with $\varepsilon(B)=A_0B^{-1/12}$,
\begin{equation}\label{eq:v93-endpoint}
 \mathcal E_M\le C_MB^{-1/192}
       +\mathfrak I_{\varepsilon(B),M}(r)
       +\mathfrak C_{\varepsilon(B),M}(r)
       +C(1+\varepsilon(B)^{-1})r.
\end{equation}
The first and last terms are estimated. The two positive physical
forcing terms are retained, not asserted to vanish. In particular,
vanishing of their sum along a legal choice $r(B)$ with
$(1+\varepsilon(B)^{-1})r(B)\to0$ is sufficient for the original
pointwise assertion. Necessity of this additional forcing condition
is not claimed.
\end{theorem}
\begin{proof}
Use \eqref{eq:v93-upper-finite} with the original full bad source.
Its backward average is at most $2A_rb^\varepsilon$. Fix
$B,r,\rho$ and every auxiliary band before the collision limit.
The fixed-width estimate gives $C\varepsilon^{1/16}$ for the
normalized average, while Theorem~\ref{thm:v93-gain} bounds the
noncritical gain by $C(1+\varepsilon^{-1})r$. Only then let
$\rho\downarrow0$ in the remaining critical term. No integral or
supremum is interchanged with this limit. Finally insert the positive
scalar remainder and use
$\varepsilon(B)^{1/16}=O(B^{-1/192})$.
For example $r(B)=c\varepsilon(B)^2$ makes the explicit last term
vanish, but does not estimate either forcing term. Every such width
is fixed before its collision limit, with $B$ outermost.
\end{proof}

The entrance measure may contain critical directional atoms; the
critical growth may concentrate after the collision normalization.
The uniform per-sheet modulus does not bound the complete transverse
sum. These are precisely the reasons that
\eqref{eq:v93-endpoint} is not yet a charge-zero theorem. It is a
pointwise source comparison, with signed physical losses and artificial
prefix losses accounted for before the upper estimate, rather than
another bounded normalization or finite-roof-norm conclusion.
